\begin{textsample}{POS Dim 6 – sc – Score 26.00 – sc/sc\_em\_45.txt}  \label{ex:f6_pos_013}
1.1 \textbf{INTRODUÇÃO} A Resolução nº 3, de 21 de novembro de 2018, do Ministério de Educação ( MEC ), atualiza as Diretrizes Curriculares Nacionais para o Ensino Médio ( DCNem ).
O documento é um referencial legal que define obrigatoriedades para o ensino médio, a partir do pressuposto de que esta etapa da educação básica é direito de todos, é dever do Estado e da família e será promovido e incentivado com a colaboração da sociedade.
Novos termos e estrutura do ensino médio são definidos visando ao pleno desenvolvimento da pessoa, ao seu preparo para o exercício da cidadania e à sua qualificação para o trabalho.
No capítulo 2, que trata do referencial legal e conceitual, especificamente no parágrafo 6º, itens II e III, tem -se a definição de dois componentes \textbf{estruturantes} do Novo Ensino Médio : formação geral básica e itinerário \textbf{formativo}.
Este último é definido como : > cada conjunto de unidades curriculares ofertadas pelas instituições e redes de ensino que possibilitam ao estudante aprofundar seus conhecimentos e se preparar para o prosseguimento de estudos ou para o mundo do trabalho de forma a contribuir para a construção de soluções de \textbf{problemas} específicos da sociedade ( Resolução 3, de 21/11/2018 ).
Para a área de Matemática, e suas \textbf{tecnologias}, de acordo com a Resolução nº 3, de 21 de novembro de 2018, os itinerários \textbf{formativos} ( \textbf{IF} ) devem enfatizar o : > [... ] aprofundamento de conhecimentos \textbf{estruturantes} para aplicação de diferentes conceitos matemáticos em contextos sociais e de trabalho, estruturando arranjos curriculares que permitam estudos em resolução de \textbf{problemas} e análises complexas, funcionais e não-lineares, análise de dados estatísticos e probabilidade, geometria e topologia, robótica, automação, inteligência artificial, programação, jogos \textbf{digitais}, sistemas dinâmicos, dentre outros, considerando o contexto local e as possibilidades de oferta pelos sistemas de ensino.
Os itinerários \textbf{formativos} assumem grande importância na formação de estudantes e têm, na Portaria nº 1.432, de 28 de dezembro de 2018, a definição de seus pressupostos e estrutura.
Constituem a parte flexível e diversificada do currículo, enfatizando o aprofundamento e a ampliação de aprendizagens em uma ou mais áreas de conhecimento e/ou na formação técnica e profissional, com carga horária total mínima de 1.200 horas.
O aprofundamento em uma área do conhecimento, ou na formação técnica profissional, tem por finalidade a formação integral do sujeito.
Integra a formação geral, possibilitando que os(as) estudantes desenvolvam seus aspectos cognitivos, com ênfase em suas aptidões.
Impactos tais como continuidade nos estudos, atuação no mundo do trabalho, resolução de demandas complexas na vida cotidiana e exercício da cidadania são esperados, quando se propõem aprofundamentos via itinerários \textbf{formativos} ( \textbf{IF} ).
De acordo com a Resolução nº 3, de 21 de novembro de 2018, e a Portaria nº 1.432, de 28 de dezembro de 2018, os Itinerários \textbf{Formativos} devem ser fundamentados por quatros eixos \textbf{estruturantes}, listados a seguir. - Investigação científica — Tem como ênfase ampliar a capacidade do estudante de investigar a realidade.
Neste movimento de investigação, ele compreende, valoriza e aplica o conhecimento sistematizado, por meio da realização de práticas e produções científicas, com temáticas de seu interesse, relacionadas a uma ou mais áreas de Conhecimento e à Formação Técnica e Profissional. - Processos criativos — O estudante desenvolverá a capacidade de idealizar e realizar projetos de forma criativa, sendo estes associados à temáticas de seu interesse, a uma ou mais áreas do conhecimento ou à educação técnica e profissional. - \textbf{Mediação} e intervenção sociocultural — A ênfase neste eixo é desenvolver a capacidade de realizar projetos que contribuam com a sociedade e o meio ambiente.
Esta deverá acontecer por meio de conhecimentos relacionados a uma ou mais áreas de conhecimento ou pela formação técnica e profissional. - \textbf{Empreendedorismo} — O destaque, neste eixo, é o desenvolvimento da capacidade de mobilizar conhecimentos em diferentes áreas.
As \textbf{tecnologias} devem assumir um papel de destaque, auxiliando os(as) estudantes em seus projetos de vida.
No Currículo Base do Território Catarinense ( CBTC ), os itinerários \textbf{formativos} ( \textbf{IF} ) na área de Matemática e suas \textbf{tecnologias} visam ao aprofundamento de conhecimentos e à sua aplicação em diversos contextos, tal como pressupõem os documentos que legalizam o ensino médio, fundamentando -se nas habilidades dos eixos \textbf{estruturantes} e em sua associação com as habilidades específicas da Matemática, e suas \textbf{tecnologias}, como se observa no quadro 1.
Quadro 1 — Habilidades gerais e específicas dos eixos \textbf{estruturantes} do itinerário de Matemática, e \textbf{tecnologias} | Eixos \textbf{estruturantes} e habilidades gerais | Habilidades dos eixos \textbf{estruturantes} | Habilidade dos eixos \textbf{estruturantes} associada à área da matemática e das \textbf{tecnologias} | |---|---|---| | Investigação científica : fazer e pensar científico | → Identificar, selecionar, \textbf{processar} e analisar dados, fatos e evidências com curiosidade, atenção, criticidade e ética, utilizando, inclusive, o apoio de \textbf{tecnologias} \textbf{digitais}. → Posicionar -se com base em critérios científicos, éticos e estéticos, utilizando dados, fatos e evidências para \textbf{respaldar} conclusões, opiniões e \textbf{argumentos} por meio de afirmações claras, \textbf{ordenadas}, coerentes e \textbf{compreensíveis}, sempre respeitando valores universais, como liberdade, democracia, justiça social, pluralidade, solidariedade e \textbf{sustentabilidade}. → Utilizar informações, conhecimentos e ideias \textbf{resultantes} de investigações científicas para criar ou propor soluções para \textbf{problemas} diversos. | → Investigar e analisar situações-problema, identificando e selecionando conhecimentos matemáticos relevantes para uma dada situação, elaborando modelos para sua representação. → Levantar e \textbf{testar} hipóteses sobre variáveis que interferem na explicação ou resolução de uma situação-problema, elaborando modelos com a linguagem matemática para analisá -la e avaliar sua adequação em termos de possíveis limitações, eficiência e possibilidades de generalização. → Selecionar e sistematizar, com base em estudos e/ou pesquisas ( bibliográficas, exploratórias, de campo, experimentais, etc. ) em fontes confiáveis, informações sobre a contribuição da Matemática na explicação de fenômenos de natureza científica, social, profissional, cultural, de processos tecnológicos, identificando os diversos pontos de vista, posicionando -se mediante argumentação, com o cuidado de citar as fontes dos recursos utilizados na pesquisa e buscando apresentar conclusões com o uso de diferentes mídias. | | Processos criativos : pensar e fazer criativo | → Reconhecer e analisar diferentes manifestações criativas, artísticas e culturais por meio de vivências presenciais e virtuais que ampliem a visão de mundo, a sensibilidade, a criticidade e a criatividade. → Questionar, modificar e adaptar ideias existentes e criar propostas, obras ou soluções criativas, originais ou inovadoras, avaliando e assumindo riscos para lidar com as incertezas e colocá -las em prática. → Difundir novas ideias, propostas, obras ou soluções por meio de diferentes linguagens, mídias e plataformas, analógicas e \textbf{digitais}, com confiança e \textbf{coragem}, assegurando que alcancem os interlocutores pretendidos. | → Reconhecer produtos e/ou processos criativos por meio de fruição, vivências e reflexão crítica na produção do conhecimento matemático e sua aplicação no desenvolvimento de processos tecnológicos diversos. → Selecionar e mobilizar intencionalmente recursos criativos relacionados à Matemática para resolver \textbf{problemas} de natureza diversa, incluindo aqueles que permitam a produção de novos conhecimentos matemáticos, comunicando com precisão suas ações e reflexões relacionadas a constatações, interpretações e \textbf{argumentos}, bem como adequando -os a situações originais. → Propor e \textbf{testar} soluções éticas, estéticas, criativas e inovadoras para \textbf{problemas} reais, considerando a aplicação de conhecimentos matemáticos associados ao domínio de operações e relações matemáticas simbólicas e formais, de modo a desenvolver novas abordagens e estratégias para enfrentar novas situações. | | \textbf{Mediação} e intervenção sociocultural : convivência e atuação sociocultural e \textbf{ambiental} | → Reconhecer e analisar questões sociais, culturais e \textbf{ambientais} diversas, identificando e incorporando valores importantes para si e para o coletivo que assegurem a tomada de decisões conscientes, \textbf{consequentes}, colaborativas e responsáveis. → Compreender e considerar a situação, a opinião e o sentimento do outro, agindo com empatia, flexibilidade e resiliência para promover o diálogo, a colaboração, a \textbf{mediação} e a resolução de conflitos, o combate ao preconceito e a valorização da diversidade. → Participar ativamente da proposição, implementação e avaliação de solução para \textbf{problemas} socioculturais e/ou \textbf{ambientais} em nível local, regional, nacional e/ou global, responsabilizando -se pela realização de ações e projetos voltados ao bem comum. | → Identificar e explicar questões socioculturais e \textbf{ambientais}, aplicando conhecimentos e habilidades matemáticas para avaliar e tomar decisões em relação ao que foi observado. → Selecionar e mobilizar intencionalmente conhecimentos e recursos matemáticos para propor ações individuais e/ou coletivas de \textbf{mediação} e intervenção sobre \textbf{problemas} socioculturais e \textbf{problemas} \textbf{ambientais}. → Propor e \textbf{testar} estratégias de \textbf{mediação} e intervenção para resolver \textbf{problemas} de natureza sociocultural e de natureza \textbf{ambiental} relacionados à Matemática. | | \textbf{Empreendedorismo} : autoconhecimento e projeto de vida | → Reconhecer e utilizar qualidades e fragilidades pessoais com confiança para superar desafios e alcançar objetivos pessoais e profissionais, agindo de forma proativa e empreendedora e \textbf{perseverando} em situações de \textbf{estresse}, frustração, fracasso e \textbf{adversidade}. → Utilizar estratégias de planejamento, organização e \textbf{empreendedorismo} para estabelecer e adaptar metas, identificar caminhos, mobilizar apoios e recursos, para realizar projetos pessoais e produtivos com foco, \textbf{persistência} e efetividade. → Refletir continuamente sobre o seu próprio desenvolvimento e sobre seus objetivos presentes e futuros, identificando aspirações e oportunidades, inclusive relacionadas ao mundo do trabalho, que orientem escolhas, esforços e ações em relação à sua vida pessoal, profissional e cidadã. | → Avaliar como oportunidades, conhecimentos e recursos relacionados à Matemática podem ser utilizados na concretização de projetos pessoais ou produtivos, considerando as diversas \textbf{tecnologias} disponíveis e os impactos socioambientais. → Selecionar e mobilizar intencionalmente conhecimentos e recursos da Matemática para desenvolver um projeto pessoal ou um empreendimento produtivo. → Desenvolver projetos pessoais ou produtivos, utilizando processos e conhecimentos matemáticos para formular propostas concretas, articuladas com o projeto de vida. | Fonte : Elaboração dos autores com base na Portaria nº 1.432, de 28/12/2018.
O itinerário \textbf{formativo} de Matemática e suas \textbf{tecnologias} visa a ampliar os conceitos estudados na formação geral por meio do desenvolvimento das habilidades apresentadas no quadro 1.
Além de explorar potenciais, o aprofundamento na área de conhecimento, permite ao jovem concluir o ensino médio com um diferencial em seu currículo.
Este é conquistado quando seu currículo escolar possibilitou maior tempo de estudo nas unidades curriculares de seu interesse, em consonância com seu projeto de vida.
Sendo assim, ao final do ensino médio, espera -se que o estudante tenha alcançado as habilidades necessárias ao desenvolvimento das dez competências gerais da BNCC.
A interdependência entre essas dez competências, as habilidades dos eixos \textbf{estruturantes} e as habilidades dos aprofundamentos na área de Matemática pode ser observada no mapa conceitual representado na \textbf{figura} 1.

% matched lemmas: adversidade, ambiental, argumento, compreensível, consequente, coragem, digital, empreendedorismo, estresse, estruturante, figura, formativo, if, introdução, mediação, ordenado, perseverar, persistência, problema, processar, respaldar, resultante, sustentabilidade, tecnologia, testar
\end{textsample}
